\section{Planteamiento del Problema}
\label{sec:problema}
El problema se presenta en cinco partes: la situación actual, a quién afecta y
por qué importa, lo que todavía no se sabe de él, su enunciado y la viabilidad
de atenderlo.

\subsection{La Situación Actual}
En muchas instituciones educativas la entrada de las personas se controla
revisando credenciales a simple vista o anotando registros a mano. Ambos
métodos dependen de la atención de quien los aplica y no siempre convencen a
quienes los usan: en una escuela de Perú, el 87.5\,\% de los 40 docentes y
administrativos encuestados no se sentía cómodo con el registro manual de
asistencia \parencite{torres2019qr}.

El TESCHA no es la excepción. Los vigilantes revisan a simple vista la
credencial física de cada estudiante. La revisión se complica en el turno
vespertino durante los 15 a 20 minutos previos a la primera clase, cuando la
mayor parte de los estudiantes llega en poco tiempo y todo depende de una o
dos personas. En ese momento se forman filas, una credencial prestada, vencida
o ajena puede pasar sin que nadie lo note y, al final del día, no queda ningún
registro de quién entró ni a qué hora.

Conviene no confundir este problema con el pase de lista en el salón. Los dos
nacen de una identificación débil, pero un registro en la entrada solo prueba
que la persona entró al plantel, no que asistió a una clase. Por eso el
control de asistencia por materia queda fuera del proyecto.

\subsection{A Quién Afecta y Por Qué Importa}
El problema afecta a tres grupos:
\begin{itemize}
	\item \textbf{Estudiantes.} Pierden tiempo en la fila, pueden llegar tarde a
	      su primera clase y conviven en un plantel donde la entrada de personas
	      ajenas no se controla con rigor.
	\item \textbf{Personal de vigilancia.} Revisa cientos de credenciales en
	      pocos minutos sin ninguna herramienta que le ayude a detectar una
	      credencial ajena, vencida o falsa.
	\item \textbf{Dirección y Departamento de Control Escolar.} No tienen datos
	      de ingreso para detectar a tiempo a los estudiantes que dejan de venir,
	      una señal temprana de posible deserción, ni para decidir cuántos accesos
	      abrir y con cuánto personal en cada horario.
\end{itemize}

Sus consecuencias son de cuatro tipos. Las \textbf{filas} retrasan el inicio
de la jornada. La teoría de colas, que ya se propuso para reducir la espera de
estudiantes en el área de control escolar de otro plantel del Tecnológico
Nacional de México \parencite{talavera2025espera}, ayuda a entender por qué:
la fila crece cuando llegan más personas por minuto de las que se alcanzan a
revisar. La \textbf{suplantación} es un riesgo de seguridad, porque permite la
entrada de personas que no pertenecen a la comunidad. La \textbf{falta de
	datos} impide tomar decisiones con evidencia. Por último, la \textbf{carga de
	trabajo} del vigilante crece justo cuando más atención se necesita.

\subsection{Lo que Todavía No Sabemos}
Hoy no existe ninguna medición del acceso del TESCHA, así que la magnitud del
problema todavía no se puede expresar con cifras. Falta conocer:
\begin{itemize}
	\item \textbf{Cómo funciona hoy el acceso:} cuántos estudiantes llegan por
	      minuto en la hora pico, cuánto tarda la revisión de cada uno, qué tan
	      larga se hace la fila y con qué frecuencia se olvidan o se prestan
	      credenciales.
	\item \textbf{Si los estudiantes pueden y quieren usar el sistema:} si tienen
	      celular propio con conexión y batería al llegar, y si estarían
	      dispuestos a mostrar un código QR para entrar.
	\item \textbf{Qué tan rápido responde el lector:} cuánto tarda en leer el
	      código y consultar al servidor, y cómo influyen la luz, el brillo de la
	      pantalla y la distancia. No se encontró literatura arbitrada que reporte
	      ese tiempo en el lector elegido, así que se medirá con el prototipo.
	\item \textbf{Qué obligaciones legales aplican:} como organismo público del
	      Estado de México, el TESCHA debe cumplir la ley estatal y la ley general
	      de protección de datos personales \parencite{edomex2017datos,
		      lgpdppso2025}, y el sistema guardará la hora de entrada de cada
	      estudiante.
\end{itemize}
Los dos primeros puntos se resuelven en el diagnóstico de octubre con
observación, encuesta y entrevistas (ver Metodología de Medición); el tercero,
en las pruebas en el salón, y el cuarto, con la revisión del aviso de
privacidad.

\subsection{Enunciado del Problema}
El acceso de los estudiantes al TESCHA carece de un mecanismo de validación
automatizado. Como consecuencia, en la hora pico del turno vespertino se
forman filas, una credencial prestada o ajena puede superar la revisión
visual y la institución no cuenta con un registro confiable de ingresos que
le sirva para tomar decisiones. De este enunciado surge la pregunta que guía
el proyecto, presentada junto con los objetivos más adelante.

\subsection{Por Qué Es Viable Atenderlo}
El uso de códigos QR para controlar el acceso físico ya se ha propuesto antes
\parencite{kao2011physical}, y su costo es bajo frente a otras opciones.
Comparado con las tarjetas NFC, el QR se procesa más rápido, aunque no trae
cifrado propio \parencite{harnaningrum2024comparison}. El proyecto compensa
esa carencia con un código que cambia cada pocos segundos y va firmado, de
modo que un código reenviado o alterado se rechaza
\parencite{subpratatsavee2015hmac}. Además, el estudiante usa su propio
celular, así que no hay que fabricar tarjetas. Comparado con el
reconocimiento facial, el QR no necesita equipo costoso ni guarda datos
biométricos, cuya aceptación entre universitarios depende mucho de la
confianza y la seguridad que perciben \parencite{rukhiran2023user}.

Atenderlo también es necesario. Sin datos no se sabe cuánto espera un
estudiante para entrar, qué tan seguido se usan credenciales ajenas ni si la
comunidad aceptaría el sistema. Medir el acceso actual y compararlo con el
prototipo permite decidir con evidencia, y no con suposiciones, si conviene
extender el sistema a los demás accesos del plantel.
